\documentclass[10pt,a4paper]{amsart}
\usepackage[margin=19mm]{geometry}
\usepackage{amsmath,amssymb,mathtools}
\usepackage[scr=rsfs,cal=euler]{mathalfa}
\usepackage{xfrac}
\usepackage[unicode,hidelinks,bookmarks=false]{hyperref}
\newcommand{\ie}{i.e.\ }
\newcommand{\cf}{cf.\ }
\newcommand{\resp}{respectively\ }
\newcommand{\Subsection}{\subsection}
\providecommand{\nobreakdash}{\nobreak\hbox{-}}
\setlength{\emergencystretch}{2em}
\multlinegap0pt
\usepackage{mathtools}
\usepackage{enumitem}
\newtheorem{theorem}{Theorem}[section]
\newtheorem{lemma}[theorem]{Lemma}
\newtheorem{corollary}[theorem]{Corollary}
\newtheorem{proposition}[theorem]{Proposition}
\newtheorem{remark}[theorem]{Remark}
\newtheorem{definition}[theorem]{Definition}
\newcommand{\diam}{\operatorname{diam}}
\newcommand{\dist}{\operatorname{dist}}
\newcommand{\ZZ}{\mathbb Z}
\newcommand{\RR}{\mathbb R}
\newcommand{\QQ}{\mathbb Q}
\newcommand{\llabs}{\ll}
\begin{document}
\title[Distances in Planar Integral Point Sets]{Distances in Planar Integral Point Sets}
\author{Jozsef Solymosi}
\date{}
\maketitle
\noindent\emph{Source and licence.} Distances in Planar Integral Point Sets. Version arXiv:2606.26311v1 (CC BY 4.0). This material is distributed under the Creative Commons Attribution 4.0 International licence, http://creativecommons.org/licenses/by/4.0/, as recorded in the arXiv OAI licence metadata for this exact version and shown on the abstract page https://arxiv.org/abs/2606.26311v1. Full rights notice: licenses/arxiv-2606.26311-CC-BY-4.0.txt.\par\smallskip
\noindent\textit{Editorial scope.} Counted unit: the original \texttt{lemma} environment \texttt{lem:normalization} at source lines 133--153 together with its complete proof at lines 155--185; the delivered document keeps source lines 101--186 so that every referenced label and every citation resolves inside it. Native editable source is the paper's own concatenated TeX; the AMS wrapper is editorial formatting only. No publisher class, upstream script or unrelated artwork is redistributed.
\section{Three anchors and bounded triple intersections}

We begin with a small divisibility fact that will be used twice.

\begin{lemma}[Integer root bound]\label{lem:root}
Let
\[
        f(T)=a_dT^d+a_{d-1}T^{d-1}+\cdots+a_0\in\ZZ[T]
\]
be a nonzero polynomial of degree at most \(D\), and suppose that
\(|a_i|\le H\) for all \(i\).  If \(R\in\ZZ\) and \(f(R)=0\), then
\[
        |R|\le H.
\]
\end{lemma}

\begin{proof}
If \(R=0\), there is nothing to prove.  Otherwise let \(j\) be the smallest
index such that \(a_j\ne0\).  Then
\[
        f(T)=T^jg(T),
        \qquad
        g(T)=a_j+a_{j+1}T+\cdots+a_dT^{d-j},
\]
with \(g(0)=a_j\ne0\).  Since \(f(R)=0\) and \(R\ne0\), we have \(g(R)=0\).
Thus
\[
        a_j=-R\bigl(a_{j+1}+a_{j+2}R+\cdots+a_dR^{d-j-1}\bigr),
\]
so \(R\mid a_j\).  Hence \(|R|\le |a_j|\le H\).
\end{proof}


\begin{lemma}[Quadratic coordinate normalization]\label{lem:normalization}
Let \(P\) be an integral point set, and let \(A,B,C\in P\) be non-collinear.
After an isometry we may write
\[
        A=(0,0),\qquad B=(d,0),\qquad
        C=\left(\frac{U}{L},\frac{\sqrt{k}V}{L}\right),
        \qquad L=2d,
\]
where \(d=|AB|\in\ZZ_{>0}\), \(U,V\in\ZZ\), \(V\ne0\), and \(k\) is a positive
squarefree integer.  Moreover every point \(X\in P\) has coordinates
\[
        X=(x,\sqrt{k}y)
        \qquad\text{with }x,y\in\QQ.
\]
If \(|AB|,|AC|,|BC|\le a\), then
\[
        L\le 2a,
        \qquad |U|\le 2a^2,
        \qquad kV^2\le 4a^4.
\]
\end{lemma}

\begin{proof}
Place \(A=(0,0)\) and \(B=(d,0)\).  If \(c=|AC|\) and \(e=|BC|\), then
\[
        C_x=\frac{d^2+c^2-e^2}{2d}.
\]
Thus \(C_x=U/L\) with \(L=2d\) and \(U\in\ZZ\).  Since
\[
        C_y^2=c^2-\frac{U^2}{L^2},
\]
we may write \(C_y=\sqrt{k}V/L\), with \(k\) squarefree and \(V\in\ZZ\).  The
non-collinearity of \(A,B,C\) gives \(V\ne0\).

Now let \(X=(x,Y)\in P\).  Since \(|XA|\) and \(|XB|\) are integers, subtracting
squared distances gives
\[
        2dx=|XA|^2+d^2-|XB|^2\in\ZZ,
\]
so \(x\in\QQ\).  Since \(|XC|\) is also integral,
\[
        2C\cdot X=|XA|^2+|AC|^2-|XC|^2\in\ZZ.
\]
The term \((U/L)x\) is rational, so \((\sqrt{k}V/L)Y\) is rational.  As
\(V\ne0\), this gives \(Y=\sqrt{k}y\) with \(y\in\QQ\).

Finally, if the triangle has longest side at most \(a\), then \(L=2d\le2a\).
Also \(|U|/L=|C_x|\le |AC|\le a\), so \(|U|\le2a^2\).  Similarly,
\[
        \frac{kV^2}{L^2}=C_y^2\le |AC|^2\le a^2,
\]
whence \(kV^2\le4a^4\).
\end{proof}


\end{document}
